\chapter{ROS 2 Controllers and Communication}
\label{ch:ros2}

\section{Nodes, graph, and interfaces}
ROS 2 nodes are independent processes or components connected by typed
interfaces. Topics stream state/events, services perform request-response
operations, and actions represent longer tasks with acceptance, feedback,
completion, and cancellation. The project uses all three general styles in
its ecosystem: topics for joint/gripper/stop streams, the
\code{FollowJointTrajectory} action for arm motion, and controller-manager
services for inspection. In the task path, the detachable-joint endpoints are
Gazebo Transport topics, not ROS services. ROS 2's own guide explains the
topic/service/action distinction \cite{ros2_interfaces}.

\section{Controller manager and hardware interfaces}
The URDF instantiates the external Gazebo simulated hardware system. The YAML
sets controller-manager update rate to \SI{100}{Hz} and simulation time. It
loads:
\begin{enumerate}
 \item \code{joint_state_broadcaster/JointStateBroadcaster};
 \item \code{joint_trajectory_controller/JointTrajectoryController} for
 all six arm joints;
 \item \code{position_controllers/JointGroupPositionController} for the
 single \code{gripper_driver} joint.
\end{enumerate}
The arm controller claims a \code{position} command interface and reads
\code{position} and \code{velocity} state interfaces. The gripper position
controller receives an array with one driver position. The five mimic joints
have state interfaces but no direct command interface in the URDF.

The joint trajectory controller exposes an action endpoint
\begin{center}\code{/arm_controller/follow_joint_trajectory}\end{center}
using \code{control_msgs/action/FollowJointTrajectory}. The project sends a
joint-name list and time-stamped \code{JointTrajectoryPoint} values with
positions and velocities. It checks acceptance, terminal status, action error
code, then compares final joint feedback with the final commanded point.
Controller package documentation is cited for interface semantics, while the
project YAML/code is authoritative for what is configured here
\cite{ros2_control_jtc}.

\section{ROS topics and message types}
\begin{longtable}{p{.39\textwidth}p{.26\textwidth}p{.29\textwidth}}
\caption{Important ROS interfaces visible in project source.}\label{tab:ros-interfaces}\\
\toprule Name & Type & Actual use\\\midrule\endfirsthead
\toprule Name & Type & Actual use\\\midrule\endhead
\code{/joint_states} & \code{sensor_msgs/msg/JointState} & Broadcaster feedback; task and manual arm node read arm positions; task also reads gripper driver angle.\\
\code{/arm_controller/follow_joint_trajectory} & \code{control_msgs/action/FollowJointTrajectory} & Task and manual arm/demo nodes send multi-point arm goals.\\
\code{/gripper_controller/commands} & \code{std_msgs/msg/Float64MultiArray} & Task and manual gripper node publish the single driver setpoint.\\
\code{/gp7/stop} & \code{std_msgs/msg/Empty} & Task node requests cancellation; manual arm node also handles stop. Not a safety-rated emergency stop.\\
\code{/gp7/goal_index} & \code{std_msgs/msg/Int32} & Manual arm preset selector (0--3), not used by assignment task node.\\
\code{/gp7/gripper_cmd} & \code{std_msgs/msg/Float64} & Manual gripper-node input, translated into controller command array.\\
\code{/clock} & \code{rosgraph_msgs/msg/Clock} & Gazebo clock bridged to ROS; nodes use simulation time.\\
\code{/controller_manager/list_controllers} & ListControllers service & Human-operated controller readiness check documented in README; task launch spawners wait for startup.\\
\bottomrule
\end{longtable}

The exact action and topic names are absolute names in the source. Controller
names in the YAML are \code{joint_state_broadcaster}, \code{arm_controller},
and \code{gripper_controller}. The Gazebo system creates the controller
manager. Both Gazebo launch descriptions enforce broadcaster startup before
or alongside the command controllers; the generic launch has a process-exit
event handler that starts arm and gripper spawners only after the broadcaster
spawner finishes.

\section{Gazebo Transport interfaces used alongside ROS}
These endpoints are native Gazebo Transport messages, despite their topic-like
names. The task code imports them from the vendor binding rather than through
\code{ros_gz_bridge}:
\begin{table}[htbp]
\centering\small
\caption{Gazebo Transport messages in \code{PickPlace}.}\label{tab:gz-interfaces}
\begin{tabularx}{\textwidth}{l l X}
\toprule Gazebo topic & Message type & Direction and purpose\\\midrule
\code{/gp7/grasp/attach} & \code{gz.msgs.Empty} & Task publishes; detachable-joint plugin attaches Box B.\\
\code{/gp7/grasp/detach} & \code{gz.msgs.Empty} & Task publishes; plugin removes attachment.\\
\code{/gp7/grasp/state} & \code{gz.msgs.StringMsg} & Plugin publishes text state such as \code{attached} / \code{detached}; task subscribes.\\
\code{/model/box_b/pose} & \code{gz.msgs.Pose_V} & SDF pose-publisher system sends model pose vector; task subscribes and selects the \code{box_b} entry.\\
\bottomrule
\end{tabularx}
\end{table}
The launch does not bridge these four topics into ROS. Only the simulation
clock has the explicit ROS--Gazebo bridge in the task launch. This split is
why a ROS topic listing alone does not show every interface used in the task.

\section{Robot-state publisher and TF}
\code{robot_state_publisher} receives the robot description and joint
states. It publishes fixed transforms for fixed URDF joints and dynamic
transforms for actuated joints. The task launch sets \code{use_sim_time}
and creates a \code{ros_gz_bridge} parameter bridge for \code{/clock}.
RViz in the separate display launch reads these transforms to render the URDF.
The assignment task itself does not rely on an RViz node.

\section{Manual control nodes versus the task node}
\code{ArmControllerNode} in \file{scripts/arm_controller_node.py} listens
for one of four preset indices, subscribes to joint states, and sends action
goals. Its stop callback cancels the active action and attempts a hold using
the last observed joint position. \code{GripperControllerNode} accepts a
simple scalar, checks finite value and $[-0.5,0.15]$ bounds, and publishes a
one-element array. \code{DemoSequenceNode} runs scripted joint waypoints and
gripper commands. These are separate manual/demo utilities; none performs
the integrated feedback-checked Box A/Box B workflow.

The task node itself bypasses these two wrapper nodes: it publishes directly
to \code{/gripper_controller/commands}, sends the arm action directly, and
subscribes to \code{/gp7/stop}. This matters when tracing nodes in a graph:
the arm and gripper wrappers need not be running during \code{pick_place.launch.py}.

\section{Communication diagram}
\Cref{fig:architecture} shows launch, state publisher, controller manager,
Gazebo, and the task. The clock path is bridged from Gazebo Transport into
ROS. Conversely, \code{pick_place_node.py} creates its own
\code{gz.transport13.Node}; it directly advertises and subscribes to grasp
and model-pose topics. Therefore not every connection in the task is a ROS 2
topic, and topic names alone do not identify the middleware type.

\section{Command, feedback, and failure handling}
An action goal has a lifecycle: the controller can reject it; an accepted
goal can finish with a result code; or the task can request cancellation. The
task's \code{move} function waits for the goal response, stores the active
handle, waits with a duration-based timeout, cancels on exception, checks
success, waits for settling, and compares live joint feedback. During carrying,
it also checks that pose and grasp-state feedback remain fresh and that the
largest observed orientation error remains within one degree. This adds
application-level checking around the controller action.
